\chapter{Le C++ : ce qu'il ajoute, et ce que Nkentseu en garde}

Le C++ ajoute au C de quoi rendre les erreurs \emph{impossibles} plutôt que
seulement improbables. Ce chapitre couvre le langage, puis dit --- en le mesurant
--- ce que ce dépôt en emploie et ce qu'il refuse. La seconde partie compte
autant que la première : apprendre le C++ d'un livre standard, c'est apprendre
des choses qu'il faudra désapprendre ici.

\section{La norme employée}

\begin{center}
\begin{tabular}{@{}lll@{}}
\toprule
\textbf{Norme} & \textbf{Projets} & \textbf{Mesure} \\
\midrule
C++17 & 46 & \texttt{cppdialect("C++17")} \\
C++20 & 2  & \texttt{cppdialect("C++20")} \\
\bottomrule
\end{tabular}
\end{center}

Écrivez du C++17. Les deux projets en C++20 ont une raison particulière ; la
règle est 17.

\section{Classes et structures}

Une \texttt{class} est une \texttt{struct} dont les membres sont privés par
défaut. C'est \emph{toute} la différence technique.

\begin{nkcode}{Exemple pedagogique}
class NkCompteur {
    public:                            // ce que le monde exterieur voit
        void Incrementer() { ++mValeur; }
        int32 Valeur() const { return mValeur; }   // `const` : ne modifie rien

    private:                           // ce qui n'appartient qu'a la classe
        int32 mValeur = 0;
};
\end{nkcode}

\begin{nkretenir}
\textbf{Convention du dépôt} : membres privés préfixés \texttt{m}
(\texttt{mValeur}), méthodes en \texttt{PascalCase}, types préfixés \texttt{Nk}.
Une méthode qui ne modifie pas l'objet est \texttt{const} --- ce n'est pas de la
politesse : sans lui, on ne peut pas l'appeler sur un objet constant, et la
moitié du code devient inaccessible.
\end{nkretenir}

\begin{nkretenir}
Dans ce dépôt, \texttt{struct} sert aux \textbf{données nues} (champs publics,
aucune invariante à protéger --- \texttt{NkPointer}, \texttt{NkLayoutInfo}) et
\texttt{class} à ce qui a un \textbf{état à défendre} (\texttt{NkCanvasApp},
\texttt{NkCanvasSplash}). Le compilateur ne fait pas cette distinction ; le
lecteur, si.
\end{nkretenir}

\section{Construire et détruire : RAII}

Un \textbf{constructeur} s'exécute à la naissance de l'objet, un
\textbf{destructeur} à sa mort --- automatiquement, même si une erreur survient
entre les deux.

\begin{nkcode}{Exemple pedagogique}
class NkFichier {
    public:
        NkFichier(const char *chemin) { mHandle = Ouvrir(chemin); }
        ~NkFichier()                  { if (mHandle) { Fermer(mHandle); } }

    private:
        void *mHandle = nullptr;
};

void Traiter() {
    NkFichier f("donnees.txt");
    // ... on travaille ...
}   // <- ici, ~NkFichier s'execute. Le fichier est ferme. Toujours.
\end{nkcode}

\begin{nkretenir}
C'est \textbf{RAII} : \emph{l'acquisition d'une ressource est son
initialisation}. On ne peut plus oublier de libérer, parce que ce n'est plus
vous qui le faites.

C'est la réponse du C++ aux deux pointeurs mortels du chapitre précédent --- le
pendant et la fuite. Elle vaut pour tout ce qui se prend et se rend : mémoire,
fichier, verrou, contexte GPU.
\end{nkretenir}

\begin{nkpiege}
\textbf{L'ordre de destruction est l'INVERSE de l'ordre de déclaration.} Les
membres déclarés en dernier meurent en premier.

Ce n'est pas un détail de style. Si un membre \texttt{A} tient un pointeur vers
un membre \texttt{B}, alors \texttt{A} doit être déclaré \emph{après} \texttt{B}
--- sinon \texttt{B} meurt d'abord et \texttt{A} pointe vers un mort pendant sa
propre destruction. Le compilateur ne dit rien.
\end{nkpiege}

\section{Références}

Une référence est un \textbf{autre nom} pour une variable existante. Contrairement
au pointeur, elle ne peut être ni nulle ni réaffectée.

\begin{nkcode}{Exemple pedagogique}
void Doubler(int32 &v) { v = v * 2; }   // pas de `*`, pas de `&` a l'appel
int32 n = 21;
Doubler(n);                              // n vaut 42

void Lire(const NkPointer &p);           // reference CONSTANTE : pas de copie,
                                         // et la fonction promet de ne rien changer
\end{nkcode}

\begin{nkretenir}
\textbf{Quand passer quoi} :

\begin{center}
\begin{tabular}{@{}lll@{}}
\toprule
\textbf{Vous voulez} & \textbf{Écrivez} & \textbf{Pourquoi} \\
\midrule
lire un petit objet & \texttt{int32 n} & copier 4 octets est gratuit \\
lire un gros objet & \texttt{const T \&} & pas de copie, pas de modification \\
modifier & \texttt{T \&} & l'appelant voit le changement \\
« peut-être rien » & \texttt{T *} & seul un pointeur peut être nul \\
\bottomrule
\end{tabular}
\end{center}

La dernière ligne est la vraie raison de garder des pointeurs : une référence ne
peut pas dire « absent ».
\end{nkretenir}

\section{Surcharge et valeurs par défaut}

\begin{nkcode}{Exemple pedagogique}
void Dessiner(const NkRect &r);                       // surcharge 1
void Dessiner(const NkRect &r, const NkColor &c);     // surcharge 2

void Effacer(const NkColor &c = NkColor::Blanc);      // valeur par defaut
\end{nkcode}

Deux fonctions peuvent porter le même nom si leurs paramètres diffèrent : le
compilateur choisit d'après les arguments.

\begin{nkpiege}
\textbf{Le type de retour ne participe pas au choix.} Deux fonctions qui ne
diffèrent que par lui ne compilent pas.

\textbf{Et une surcharge manquante ne se voit pas par son nom.} Si vous cherchez
\texttt{begin()} et trouvez la version \texttt{const}, vous conclurez qu'elle
existe --- alors que la version non-\texttt{const} manque. Seule l'\emph{édition
de liens} tranche : le code compile des deux côtés.
\end{nkpiege}

\section{Héritage et polymorphisme}

\begin{nkcode}{Exemple pedagogique -- le patron du depot}
class NkChapitre {                       // la classe de BASE
    public:
        virtual ~NkChapitre() = default;         // destructeur VIRTUEL
        virtual void Dessiner() = 0;             // = 0 : PUREMENT virtuelle
        virtual void Avancer(float32 dt) { }     // virtuelle, avec un defaut
};

class NkChapitre1 : public NkChapitre {  // la classe DERIVEE
    public:
        void Dessiner() override { /* ... */ }   // `override` : on REDEFINIT
};
\end{nkcode}

\texttt{virtual} veut dire : « si l'objet réel est d'une classe dérivée, appelle
\emph{sa} version ». C'est le polymorphisme --- un seul appel, plusieurs
comportements.

\begin{nkpiege}
\textbf{Un destructeur de classe de base doit être \texttt{virtual}.} Sans lui,
détruire un objet dérivé à travers un pointeur de base n'appelle que le
destructeur de la base : les membres de la dérivée ne sont jamais libérés. Aucun
avertissement, une fuite silencieuse.

\textbf{Écrivez \texttt{override}.} Sans ce mot, une faute de frappe dans le nom
ou une différence de signature crée une \emph{nouvelle} méthode au lieu d'en
redéfinir une. Le code compile, et votre méthode n'est jamais appelée. Avec
\texttt{override}, le compilateur refuse.
\end{nkpiege}

\begin{nkpiege}
\textbf{\texttt{final} interdit la redéfinition, et ce n'est pas un détail.} Dans
ce dépôt, \texttt{NkCanvasGuiApp::OnRender} est \texttt{final} : une classe qui
en dérive \emph{ne peut pas} dessiner autrement qu'avec la liste d'affichage
NKGui. Ce seul mot décide de quelle coquille vous devez partir.
\end{nkpiege}

\section{Templates}

Un \emph{modèle} écrit du code une fois pour plusieurs types.

\begin{nkcode}{Exemple pedagogique}
template <typename T>
T Maximum(const T &a, const T &b) { return (a > b) ? a : b; }

int32 x = Maximum(3, 7);          // T = int32, deduit
float32 y = Maximum(1.f, 2.f);    // T = float32
\end{nkcode}

C'est ainsi qu'est écrit \texttt{NkVector<T>}, et c'est aussi la forme de
\texttt{NkCanvasApp::Run<T>} : la coquille ne connaît pas votre classe, elle
la reçoit en paramètre de modèle.

\begin{nkpiege}
\textbf{Les erreurs de template sont longues et intimidantes.} Lisez la
\emph{première} ligne du message : elle nomme la vraie cause. Les vingt suivantes
ne font que dérouler la pile d'instanciation.

\textbf{Et le code d'un template vit dans l'en-tête}, pas dans un \texttt{.cpp} :
le compilateur doit le voir pour l'instancier avec votre type. Le mettre dans un
\texttt{.cpp} donne une erreur d'édition de liens.
\end{nkpiege}

\section{Espaces de noms}

\begin{nkcode}{Exemple pedagogique}
namespace nkentseu {
    namespace renderer {
        class NkCanvasApp { ... };
    }
}

nkentseu::renderer::NkCanvasApp app;   // nom complet
using namespace nkentseu::renderer;    // ou : on ouvre l'espace
NkCanvasApp app;
\end{nkcode}

\begin{nkpiege}
\textbf{Deux types peuvent porter le même nom dans deux espaces différents.} Ce
dépôt en a un cas réel : \texttt{NkShaderStage} existe dans
\texttt{nkentseu} \emph{et} dans \texttt{nkentseu::renderer}. Non qualifié, le
second gagne --- et l'erreur ne sort pas dans le fichier fautif.

Dans un \textbf{en-tête}, qualifiez toujours complètement. Un
\texttt{using namespace} en en-tête l'impose à tous ceux qui l'incluent.
\end{nkpiege}

\section{Ce que Nkentseu emploie, et ce qu'il refuse}

\subsection{La règle : zéro bibliothèque standard}

Nkentseu écrit ses propres conteneurs. \texttt{NkVector} au lieu de
\texttt{std::vector}, \texttt{NkString} au lieu de \texttt{std::string},
\texttt{NkAlloc}/\texttt{NkFree} au lieu de \texttt{new}/\texttt{delete}.

\begin{center}
\begin{tabular}{@{}ll@{}}
\toprule
\textbf{Ailleurs} & \textbf{Ici} \\
\midrule
\texttt{std::vector<T>} & \texttt{NkVector<T>} \\
\texttt{std::string} & \texttt{NkString} \\
\texttt{new} / \texttt{delete} & \texttt{memory::NkAlloc} / \texttt{NkFree} \\
\texttt{std::optional<T>} & \texttt{NkOptional<T>} \\
\texttt{std::cout} & \texttt{logger.Info(\dots)} \\
\texttt{printf} & \texttt{logger.Infof(\dots)} \\
\bottomrule
\end{tabular}
\end{center}

\begin{nkretenir}
La règle est \textbf{positive}, pas seulement négative : on n'écrit pas
\texttt{std::}, et \emph{si la primitive manque, on l'écrit au niveau où elle
doit vivre} --- dans Foundation, pas chez soi. Un utilitaire écrit localement
devient le troisième exemplaire d'une chose qui aurait dû être partagée.
\end{nkretenir}

\subsection{Et l'état réel, mesuré}

\begin{nkpiege}[Ce qu'un \texttt{grep} vous montrera, et comment le lire]
Le 2 septembre 2026, sur \texttt{Kernel/Foundation} et \texttt{Kernel/Runtime} :
\textbf{1\,887} occurrences de \texttt{std::}, dont \texttt{std::string} (191),
\texttt{std::vector} (160), \texttt{std::cout} (78).

\textbf{110 fichiers} emploient \texttt{std::vector} ou \texttt{std::string}, et
\textbf{trois seulement} sont des tests ou des démos. Ce n'est donc pas du code
jetable : c'est du code livré.

Ils se concentrent là où le moteur dialogue avec des bibliothèques tierces qui,
elles, parlent STL --- \texttt{NKSL} (glslang), \texttt{NKRHI} (Vulkan),
\texttt{NKRenderer}.

\textbf{La règle est donc un cap, pas un état.} On la dit franchement pour deux
raisons : parce qu'un étudiant qui grep le découvrira de toute façon, et parce
qu'une règle présentée comme acquise alors qu'elle ne l'est pas cesse d'être
respectée. Dans le code \emph{que vous écrirez} --- une application, un jeu ---
elle tient sans exception.
\end{nkpiege}

\subsection{Deux mots qui trompent}

\begin{nkpiege}
\texttt{std::move} et \texttt{std::forward} \textbf{n'allouent rien et ne
déplacent rien} : ce sont des conversions de type. Les compter comme des
« corruptions mémoire » ou des « allocations » est faux. Ce qu'ils signalent,
c'est un fichier qui raisonne encore en STL --- une surface d'exposition, pas un
défaut.
\end{nkpiege}

\begin{nkpiege}
\textbf{Une enveloppe bien nommée endort la méfiance.} \texttt{env::GetEnvVar()}
rend un \texttt{const char *} qui vaut \texttt{nullptr} quand la variable
n'existe pas --- exactement comme \texttt{getenv}. Mais son nom \emph{maison}
suggère un objet de la maison, et l'on écrit alors
\texttt{NkString v = env::GetEnvVar("X");} sans y penser. Le dépôt a failli
planter toutes ses courses GPU sur cette ligne.

\textbf{Lisez la signature de ce que vous appelez}, surtout quand vous remplacez
une fonction par une autre sur consigne.
\end{nkpiege}

\section{Ce qu'on ne trouve presque pas ici}

\begin{center}
\begin{tabular}{@{}ll@{}}
\toprule
\textbf{Construction} & \textbf{Statut dans le dépôt} \\
\midrule
exceptions (\texttt{throw} / \texttt{catch}) & rares, aux frontières tierces \\
\texttt{dynamic\_cast}, \texttt{typeid} & rares ; on préfère un champ de type \\
héritage multiple & évité \\
surcharge d'opérateurs & mesurée : \texttt{==}, \texttt{+}, \texttt{[]} sur les types mathématiques \\
\bottomrule
\end{tabular}
\end{center}

\begin{nkretenir}
Ce ne sont pas des interdits moraux. Un moteur doit tourner sur des plateformes
où les exceptions coûtent cher ou sont désactivées, et une hiérarchie profonde
rend un défaut d'entrée impossible à situer. Quand vous hésitez : regardez
comment le module voisin fait, et faites pareil.
\end{nkretenir}

\begin{nkpratique}
Écrivez une classe \texttt{NkTampon} qui alloue \texttt{n} octets par
\texttt{NkAlloc} dans son constructeur et les rend par \texttt{NkFree} dans son
destructeur. Donnez-lui une méthode \texttt{Taille() const}.

Puis provoquez les deux fautes du chapitre, et regardez ce que dit le
compilateur :

\begin{enumerate}
\item retirez le \texttt{const} de \texttt{Taille}, et appelez-la sur un
      \texttt{const NkTampon \&} ;
\item faites dériver une classe de \texttt{NkTampon}, retirez le \texttt{virtual}
      du destructeur, détruisez l'objet dérivé à travers un pointeur de base.
      \textbf{Le compilateur ne dira rien} --- c'est tout le problème. Ajoutez une
      trace dans chaque destructeur pour voir lequel manque.
\end{enumerate}
\end{nkpratique}

\begin{nkbilan}
Le C++ ajoute au C de quoi lier une donnée à ce qui la gère. \textbf{RAII} est
la plus importante de ces additions : une ressource acquise à la construction et
rendue à la destruction ne peut plus fuir, même si l'on sort par un chemin qu'on
n'avait pas prévu. Les \textbf{références} évitent la copie sans les dangers du
pointeur ; le \textbf{polymorphisme} permet d'écrire du code qui ignore le type
exact qu'il manipule --- au prix d'un destructeur virtuel qu'il ne faut jamais
oublier. Et Nkentseu \emph{écarte} la bibliothèque standard : les chapitres
suivants montrent par quoi.
\end{nkbilan}

\begin{nkquestions}
\begin{enumerate}
\item Qu'est-ce que RAII, et quel problème du C résout-il ?
\item Pourquoi un destructeur de classe de base destinée à l'héritage
      doit-il être \texttt{virtual} ?
\item Quelle différence entre une référence et un pointeur, du point de vue de
      ce qui est \emph{garanti} ?
\item Que signifie « zéro bibliothèque standard » dans ce dépôt, et l'interdit
      couvre-t-il les fonctions C ?
\item Citez deux mots du C++ dont le nom trompe sur ce qu'ils font.
\end{enumerate}
\end{nkquestions}

\begin{nkdemo}
Prouvez que RAII libère même quand on sort par le milieu.

Écrivez une classe qui affiche un message dans son constructeur et un autre dans
son destructeur. Employez-la dans une fonction qui contient trois
\texttt{return} à des endroits différents, puis appelez la fonction de façon à
emprunter chacun des trois chemins.

\textbf{Ce que la démonstration établit} : le message de destruction apparaît
\emph{trois fois sur trois}, sans qu'aucune ligne de libération ait été écrite.
Refaites ensuite la même chose avec un \texttt{free} manuel placé avant le
premier \texttt{return} seulement --- et comptez ce qui manque.
\end{nkdemo}

\begin{nklogique}
Une classe de base a un destructeur \emph{non} virtuel. On détruit un objet
dérivé à travers un pointeur de base.

\begin{enumerate}
\item Décrivez ce qui se passe exactement.
\item Le programme plante-t-il ? Justifiez --- et expliquez pourquoi la réponse
      rend ce défaut particulièrement coûteux.
\item Donnez deux façons de rendre l'erreur impossible : l'une par le
      destructeur, l'autre par la conception de la hiérarchie.
\end{enumerate}
\end{nklogique}
